% Dokumen kerja untuk pembaruan kerangka Bab 3 proposal.
\documentclass[12pt,a4paper]{report}
\usepackage[margin=3cm]{geometry}
\renewcommand{\chaptername}{Bab}
\setcounter{chapter}{2}
\begin{document}
\chapter{Metode dan Rancangan Penelitian}

\section{Deskripsi Umum Penelitian}
Uraikan alur operator mulai dari membuka profil domba, mencatat kegiatan
pemeliharaan, memindai tag RFID, hingga menelusuri bobot pada aplikasi.
Siapkan diagram arsitektur yang menunjukkan hubungan antara timbangan berbasis
Internet of Things (IoT), penyimpanan lokal, layanan penerima unggahan,
basis data, dan aplikasi seluler. Jelaskan pembagian kerja Hafiz pada
aplikasi, basis data, dan rekonsiliasi unggahan serta Lutfi pada perangkat
keras dan perangkat tegar. Keduanya menyepakati format pertukaran data dan
menjalankan pengujian integrasi.

\section{Metodologi Penelitian}
Awali bagian ini dengan bagan alir yang menunjukkan tahapan di bawah,
penanggung jawab, serta kegiatan yang telah dan akan dilakukan.

\subsection{Analisis Kebutuhan Sistem}
Turunkan kebutuhan yang akan dirancang dan diuji dari empat masalah pada
Bab~1. Untuk setiap kebutuhan, tentukan tindakan operator atau perangkat
dan hasil yang dapat diperiksa.
\begin{enumerate}
\item Operator perlu mengelola profil domba serta menyimpan dan menemukan
kembali catatan kesehatan, pemberian pakan, dan reproduksi pada profil terkait.
\item Hasil timbang yang memenuhi syarat perlu tercatat pada profil domba
yang sesuai. Operator perlu dapat menelusuri hasil tersebut melalui riwayat
bobot pada aplikasi.
\item Operator perlu mengganti tag RFID pada profil yang sama. Identitas tetap
domba dan riwayat penugasan tag perlu menghubungkan bobot yang sudah tercatat
sebelum penggantian dengan hasil penimbangan menggunakan tag baru.
\item Sistem perlu menyimpan lintasan yang selesai diunggah setelah penugasan
tag lama berakhir, beserta tag asal dan waktu pengukurannya. Pengiriman ulang
lintasan yang sama perlu dikenali agar riwayat bobot resmi bebas dari duplikasi.
\end{enumerate}

\textbf{Cakupan minimum prototipe yang akan dikembangkan dan diuji}
mencakup pengelolaan profil dan catatan pemeliharaan, pencatatan hasil
timbang IoT yang memenuhi syarat pada profil yang tepat, penggantian tag RFID
dengan riwayat bobot yang tetap dapat ditelusuri pada profil yang sama, serta
rekonsiliasi unggahan tertunda dan berulang. Setelah penggantian tag,
lintasan tertunda dari tag lama tetap tersimpan untuk penelusuran di luar
riwayat bobot resmi. Pengiriman ulang data yang sama tidak menambah catatan
bobot resmi.

\subsection{Perancangan Aplikasi dan Basis Data}
Rancang alur pengelolaan profil, pengisian catatan pemeliharaan, penggantian
tag, dan penelusuran riwayat bobot. Siapkan diagram relasi entitas yang
menunjukkan hubungan antara identitas tetap domba, penugasan tag beserta masa
berlakunya, catatan pemeliharaan, lintasan, dan catatan bobot resmi.

\subsection{Perancangan Integrasi Timbangan dan Rekonsiliasi}
Rancang format data lintasan, penomoran sampel, pemeriksaan kelengkapan,
serta konfirmasi penyimpanan kepada perangkat. Jelaskan syarat pencatatan
bobot resmi berdasarkan keabsahan hasil ukur, kelengkapan data, waktu
pengukuran yang tepercaya, dan satu penugasan tag yang cocok saat pengukuran
serta masih aktif. Siapkan diagram urutan yang memperlihatkan waktu
pengukuran, penggantian tag, penyelesaian unggahan, dan pengiriman ulang.

\subsection{Pengembangan dan Integrasi Prototipe}
Uraikan urutan pengembangan aplikasi dan layanan sisi server, penyambungan
ESP32 ke layanan penerima unggahan, serta penyajian hasil pada profil domba.
Keluaran tahap ini berupa prototipe terintegrasi yang dapat digunakan operator
untuk mencatat kegiatan pemeliharaan dan menelusuri hasil penimbangan.

\subsection{Pengujian dan Analisis Data}
Rencanakan pengujian menyeluruh dari perangkat IoT hingga aplikasi seluler.
Untuk setiap skenario, tetapkan profil dan penugasan tag awal, tindakan uji,
serta hasil yang diharapkan. Cocokkan identitas lintasan, tag, waktu ukur,
dan bobot pada log perangkat dengan rekaman basis data dan riwayat yang
benar-benar dilihat operator pada aplikasi.
\begin{enumerate}
\item \textbf{Profil dan catatan pemeliharaan.} Simpan catatan kesehatan,
pemberian pakan, dan reproduksi, lalu buka kembali profil untuk memeriksa
isi dan identitas domba. Laporkan jumlah skenario yang lulus atas seluruh
skenario yang dijalankan.
\item \textbf{Penimbangan hingga aplikasi.} Ikuti satu lintasan sejak
pemindaian RFID dan deteksi bobot stabil pada ESP32 hingga satu bobot resmi
muncul pada profil yang tepat. Bandingkan hasil timbang dengan timbangan
acuan. Hitung galat absolut dalam kilogram dan proporsi lintasan yang
menghasilkan bobot diterima. Sertakan pembacaan tidak stabil dan pembacaan
yang tampak stabil saat domba baru berada sebagian di atas platform.
\item \textbf{Penggantian tag.} Mulai dengan satu bobot resmi dari tag A.
Ganti tag A dengan tag B pada profil yang sama, lalu timbang kembali.
Periksa masa penugasan kedua tag dan pastikan operator dapat melihat bobot
lama serta bobot valid dari tag B melalui profil yang sama.
\item \textbf{Unggahan tertunda dan berulang.} Timbang domba dengan tag A
saat koneksi terputus. Setelah koneksi pulih, simpan penggantian ke tag B
sebelum unggahan lintasan A selesai. Timbang dengan tag B, selesaikan
unggahan lintasan A, lalu kirim ulang lintasan tersebut. Periksa bahwa
lintasan A tetap dapat ditelusuri dengan tag dan waktu ukur asal, bobot
resmi hanya bertambah satu dari lintasan B, dan pengiriman ulang menghasilkan
nol duplikasi. Uji pula unggahan sebagian, konflik penugasan tag, serta
waktu pengukuran yang meragukan.
\end{enumerate}
Jumlah domba, jumlah pengulangan, prosedur kalibrasi, aturan deteksi bobot
stabil, dan ambang galat masih perlu disepakati sebelum pengujian.
\end{document}
